\begin{textsample}{POS Dim 2 – human – Score 13.00 – es/es\_ef\_anos\_iniciais\_ciencias\_da\_natureza\_1\_p8.txt}  \label{ex:f2_pos_001}
A \textbf{partir} dos anos 90, o ensino de Ciências da Natureza torna -se \textbf{tema} para diversas pesquisas, que mostram que seus princípios tradicionais \textbf{levam} ao fracasso da apropriação do conhecimento e dificulta a interação entre o ensino das Ciências da Natureza e a realidade do estudante ( KRASILCHIK, 1992 ).
Com a publicação dos Parâmetros Curriculares Nacionais ( PCN ) ( BRASIL, 1997 ), no final dessa década, \textbf{buscou} -se a elaboração de \textbf{novas} propostas curriculares, que integrassem e contextualizassem conhecimentos por \textbf{meio} de \textbf{temas} transversais a todas as áreas, além de introduzir os conceitos teóricos vigentes à época, quanto ao \textbf{desenvolvimento} de competências e habilidades necessárias aos estudantes ( ESPÍRITO \textbf{SANTO}, 2009 ).
Firmando -se numa perspectiva sociocultural do ensino de Ciências, concebeu -se o conhecimento científico como uma produção sociocultural histórica que, como qualquer outra produção \textbf{humana}, contribui para o \textbf{desenvolvimento} das \textbf{capacidades} cognitivas e afetivas propriamente \textbf{humanas}.
Tal \textbf{desenvolvimento} se recria na interação dialética entre o \textbf{desenvolvimento} cultural ( história pessoal ) e o \textbf{desenvolvimento} social do sujeito ( história em \textbf{sociedade} ) ( BRASIL, 1997 ).

% matched lemmas: buscar, capacidade, desenvolvimento, humano, levar, meio, novo, partir, santo, sociedade, tema
\end{textsample}
